\cotefile{84}
\scanpage{162}
\[
x\longmapsto\pi'(x)=\pi(x)-(\chi-1)x
\]
\diagram{\begin{tikzcd}[ampersand replacement=\&,column sep=small]
0\arrow[r]\&L\arrow[r,dash]\&E\arrow[r,dash]\&L'\arrow[r,dash]\&0
\end{tikzcd}}
\deleted{\unsure{pour} \illegible{1}}.
\[
\begin{aligned}
&E\longrightarrow E\\
\deleted{\pi'(L)}\quad&(\pi'|L)=\deleted{\unsure{\chi}}\mathrm{id}_L
\qquad\Longleftrightarrow\quad\pi|L=\chi\,\mathrm{id}_L\\
&(\pi'|E/L=L')=-(\chi-1)\,\mathrm{id}_{L'}\\
&\hspace{4em}\Longleftrightarrow\quad\pi:E\longrightarrow L
\end{aligned}
\]
\[
\deleted{\unsure{\pi(x)|\pi'(x)}}\qquad
\begin{aligned}
\unsure{u|x}_1\otimes x_2+
\underbrace{\pi'_1(x_1)}_{u_1}\otimes\pi'_2(x_2)
&=u_1\otimes\underbrace{(x_2+\pi'_2(x_2))}_{x_2+\pi_2(x_2)-(\chi-1)x_2}\\
&\hspace{4em}=\underline{(2-\chi)x_2+\pi_2(x_2)}
\end{aligned}
\]
\diagram{\begin{tikzcd}[ampersand replacement=\&,row sep=small,column sep=small]
L_1\&{\subset}\& E_1\arrow[ll,bend right=25,"{\pi_1\quad\chi\text{-}\unsure{\text{triv.}|\text{tordu}}}"']\\
L_2\&{\subset}\& E_2\arrow[ll,bend left=25,"\pi_2"]
\end{tikzcd}}
\[
\Longrightarrow
\]
\diagram{\begin{tikzcd}[ampersand replacement=\&,row sep=small,column sep=tiny]
E_1\otimes E_2\arrow[rr,very thick,"?"]\arrow[dr]\arrow[dd]\&\&E_1\otimes E_2\arrow[dd]\\
\&E\arrow[ur]\&\\
E_1\otimes\deleted{\unsure{L}}E_2\underset{L_1\otimes L_2}{\oplus}L_1\otimes E_{\unsure{1|2}}\arrow[rr,"{\pi_1,\pi_2}"']\arrow[dr]\&\&L'_1\otimes L_2\\
\&L'_1\otimes L'_2\arrow[ur,dash,shorten <=7pt,shorten >=15pt]\&
\end{tikzcd}}
\diagram{\begin{tikzcd}[ampersand replacement=\&,row sep=large,column sep=small]
\unsure{0}\arrow[r,dash]\&E_1\otimes_0 E_2\arrow[r,dash]\arrow[d,"\mathrm{Tr}" ]\& E_1\otimes E_2\arrow[r,dash]\arrow[d,dash]\& L'_1\otimes L'_2\arrow[r,dash]\arrow[d,Rightarrow]\&0\\
\unsure{0}\arrow[r,dash]\&\deleted{\unsure{E|L}}L_1\otimes L_2\arrow[r,dash]\arrow[d,hook,"\text{inc. can.}"]\& E\arrow[r,two heads]\arrow[d,"?"]\&L'_1\otimes L'_2\arrow[r,dash]\arrow[d,"\chi\,\mathrm{id}"]\&0\\
0\arrow[r]\&E_1\otimes_0 E_2\arrow[r,dash]\&E_1\otimes E_2\arrow[r]\&L'_1\otimes L'_2\arrow[r,dash]\&{}
\end{tikzcd}}
\[
\begin{aligned}
&\deleted{x_1\quad\pi_1}\qquad\deleted{\alpha+}\\
&x_1\quad\pi_1(x_1)\qquad x_2,\pi_2(x_2)\\
&\alpha\; x_1\otimes x_2+\beta\,x_{\unsure{1|2}}\otimes\pi_2(x_{2})
+\gamma_{\unsure{2|{}}}\,\pi_1(x_1)\otimes x_2\unsure{+|\cdot}\delta\,\pi_1(x_1)\otimes\pi_2(x_{\unsure{2|1}})
\end{aligned}
\]
% archive sheet number: 161
\scanpage{163}
\mspage{488}
que j'avais parlé à contre temps, faute d'écoute et de discernement, avant que cette propension invétérée commence à me passer. Si elle est devenue beaucoup moins envahissante, et qu'elle a même quasiment disparu\unsure{\handwritten{\deleted{e}}|e\handwritten{´}}, ce n'est nullement comme résultat de quelque discipline que je me serais imposé\handwritten{e} (style: tu n'ouvriras le bec que si...). C'est simplement parce que l'envie m'a passé de parler, en de moments où je sens bien que c'est inutile, que cela n'apporte rien à autui ni à moi - rien du moins qui ait une valeur à mes yeux. Si à présent j'arrive souvent à sentir de telles choses, c'est sans doute que je suis devenu plus attentif. Cela aussi n'est pas venu comme fruit d'une discipline (``tu vas faire gaffe d'ouvrir grandes tes oreilles quand...''), mais je ne saurais dire comment. En tous cas, je m'en sens mieux, la vie est d'autant plus intéressante (et surtout moins bruyante!). Et les autres aussi s'en sentent mieux...

Je crois que j'ai commencé à vraiment moins parler, à partir du moment où à disparu (autant dire) cette force en moi qui me pousse à toujours vouloir rectifier ce qui m'apparaît (à tort ou à raison) comme des ``erreurs'' chez autrui - comme s'il n'était pas suffisant que je décele et rectifie les miennes! C'est la force aussi qui me poussait (et parfois encore me pousse) à vouloir à toute fin convaincre autrui de ceci ou de celà, au lieu de regarder simplement pourquoi untel préfère mordicus croire ceci plutôt que cela (qui me paraît ``ça'', et dont je voudrais bien le convaincre!); ou pourquoi je tiens tellement à ce qu'il croie cela, plutôt que ceci. Cette force quasiment universelle en nous, qui nous pousse constamment à chercher dans l'approbation d'autrui (et ne fut-ce qu'un seul...) la confirmation du bien-fondé de ce que nous tenons pour vrai - cette force profondémment enracinée dans l'égo a fini, je crois, par lâcher prise en moi. Cela a été un grand soulagement, la fin d'une dispersion d'énergie faramineuse. C'est quand j'ai fini par me rendre compte, il y a deux ans, de la portée de cette force dans ma vie, de sa nature, et de d'extraordinaire dispersion d'énergie qu'elle représentait, qu'elle s'est trouvée désamorcée - et que je me suis trouvé allégé du coup ``d'un poids de cent tonnes''. Prendre connaissance sans réticence de d'écho qu'autrui nous renvoie de notre personne, sans être lié par un désir ou ``besoin'' (si caché soit-il) d'approbation ou de confirmation - c'est cela vraiment, être ``libre de lui''. C'est un tel besoin ou désir qui constitue véritablement le ``crochet'', discret et d'une solidité à toute épreuve, par où le conflit peut ``accrocher'' en nous, et par où nous sommes (que nous le voulions ou le reconnaissions, ou non) sous la dépendance d'autrui, de
% archive sheet number: 162
\scanpage{164}
\diagram{$\begin{array}{c|c|c|}
\pi'_i&L_i&M_i\\ \hline
L_i&\chi&\tau_i\\ \hline
M_i&0&0\\ \hline
\end{array}\qquad
\begin{array}{c|c|c}
\unsure{\overline{\pi}'_i|\varpi'_i|\overline{\omega}'_i}&L_i&M_i\\ \hline
L_i&0&-\tau_i\\ \hline
M_i&0&\chi
\end{array}$}
\[
\begin{aligned}
&l_1\otimes m_2\\
&(\deleted{\unsure{l_1+m_1|l_1+m_2}}
\end{aligned}
\]
\diagram{$\begin{array}{c|c}
\phantom{1}&\phantom{11111}\\ \hline
\vphantom{\dfrac11}&
\end{array}$}
\[
\deleted{\pi'_1\otimes\mathrm{id}_{E_2}+\mathrm{id}_{E_1}\otimes}
\]
\diagram{$\begin{array}{c|ccc}
1&\deleted{\unsure{\mathrm{id}}}1&1&1\\ \hline
\pi'_1&\chi&\chi&\tau_1\otimes\mathrm{id}_{L_2}\\
\pi'_2&\chi&\mathrm{id}_{L_1}\otimes\tau_2&\chi\\
\pi'_1\otimes\pi'_2&\chi^2&\chi\,\mathrm{id}_{L_1}\otimes\tau_2&\chi\,\tau_1\otimes\mathrm{id}_{L_2}
\end{array}$}
\[
\begin{aligned}
&\alpha+\beta_1\pi'_1+\beta_2\pi'_2+\gamma\pi'_1\otimes\pi'_2\\
&\deleted{\alpha\pi'_1+\beta\pi'}
\end{aligned}
\]
\[
\deleted{\unsure{\chi}}\quad\boxed{\alpha+(\beta_1+\beta_2)\chi+\gamma\chi^2=\chi}
\]
\[
\begin{aligned}
\alpha+\beta_1\chi+
\underbrace{\beta_2\,\mathrm{id}\otimes\tau_2+\gamma\chi\,\mathrm{id}\otimes\tau_2}_{(\beta_2+\gamma\chi)\,\mathrm{id}\otimes\tau_2}
&=\mathrm{id}_{L_1}\otimes\tau_2\\
\text{i.e.}\qquad \mathrm{id}_{L_1}\otimes(\inserted{\alpha+}(\beta_1\chi\unsure{)|{}}\,\mathrm{id}_{M_2}
&+(\beta_2+\gamma\chi)\tau_2)=\\
&\mathrm{Id}\otimes\tau_2
\end{aligned}
\]
\[
\boxed{\begin{aligned}
\alpha+\beta_1\chi&=0\\
\beta_2+\gamma\chi&=1
\end{aligned}}
\]
% archive sheet number: 163
\scanpage{165}
\mspage{477}
(\textsuperscript{108}2) (6 octobre) Dans ce ``déployement de force'' il n'y a nulle intention ``agressive'' au sens courant du terme, consciente ni inconsciente, seulement un désir inconscient d'impressionner, de forcer l'estime. Il est vrai que ce terme ``forcer l'estime'' qui me revient spontanément, porte déjà une connotation de c\,o\,n\,t\,r\,a\,i\,n\,t\,e, voisine de celle ``d'agression''. Cette intention inconsciente de contrainte, perçue au niveau inconscient également, doit souvent être vécu\unsure{e} comme une sorte d'agression (alors même que ce vécu reste occulté, tout comme le\unsure{s} réactions antagonistes qu'il déclenche). En même temps, il doit se faire souvent l'amalgame de ce vécu avec des vécus analogues, remontant à l'enfance, avec le père comme protagoniste, et où celui-ci apparaît comme le principal détenteur d'autorité répressive, voire comme un rival écrasant, envié et détesté.

Même sans un tel amalgame, et indépendamment aussi de toute perception che\unsure{z} autrui d'une intention de ``contrainte'' en moi, il doit y avoir souvent la perception d'un fort d\,é\,s\,é\,q\,u\,i\,l\,i\,b\,r\,e, d'une disharmonie foncière, dans ce ``déployement de force'' exclusivement yang (dans son esprit et intention, au moins). Cette démesure est néfaste au principal intéressé, savoir moi-même, et à la limite bel et bien ``dangereux'' pour sa survie physique même (comme des incidents-santé de ces dernières années me l'ont bien montré! ). C'est cela sans doute qui a été en filigrane dans ma pensée, quand j'ecrivais qu' ``un tel déployeme\unsure{nt} de force'' était ressenti ``en tout cas comme dangereux'' - dangereux ``par nature'', un exemple donc à surtout ne pas suivre...! Un tel ressenti est sûrement suffisant pour susciter des ``réactions de défense'', en l'absence même de toute agression ou intention d'agresser.

Il est vrai que de telles relations d'ambiguité se sont reproduites après 1976, avec certains de mes élèves notamment, en des moments où tout investissement mathématique était absent, et où il n'y avait aucun ``dépoyement de force'' apparent dans ma vie. Il est vrai aussi que les `` déployements'' en question du p\,a\,s\,s\,é ont créé une r\,é\,p\,u\,t\,a\,t\,i\,o\,n, qui continue à me coller à la peau, surtout dans ma vie professionelle, et qui dans une certaine mesure se substitue à la perception de celui que je suis d\,a\,n\,s\quad l\,e\quad p\,r\,é\,s\,e\,n\,t. De plus, j'ai acquis dans le commerce de certains thèmes mathématiques une aisance telle que, même en dehors de mes périodes mathématiques et ma réputation aidant, cette aisance ou maîtrise naturelle peuvent avoir déjà l'effet du ``déployement de force'', sur des élèves peu motivés, et me faire ressentir par eux (en dépit de certains traits avenants voire rassurants) comme une sorte de Superman (un peu Superpère sur les bords!).\margin{\unsure{-}}
% archive sheet number: 164
\scanpage{166}
Pour définir un hom $F$, il faut \unsure{définir}
\diagram{\begin{tikzcd}[ampersand replacement=\&,row sep=large,column sep=small]
\& (E_1\otimes L_2)\unsure{\oplus|\amalg}(L_1\otimes E_2)=\Sigma
\arrow[dl,"{\mathrm{id}_1\otimes\alpha_2,\ \alpha_1\otimes\mathrm{id}_2}"']
\arrow[dr,"{\pi_1\otimes\mathrm{id}_{\unsure{2|L_2}},\ \mathrm{id}_1\otimes\pi_2}"]\&\\
E_1\otimes E_2\arrow[d,"F_{11}"']\&\&L_1\otimes L_2\arrow[d,"F_{22}"]\arrow[dll,"F_{21}" near start]\\
E_1\otimes\deleted{\unsure{E_2\otimes}}E_2
\arrow[dr,"{\mathrm{id}_1\otimes\psi_2,\ \psi_1\otimes\mathrm{id}_2}"']\&\&
M_1\otimes M_2\arrow[dl,"{\unsure{\varpi|\overline{\omega}}_1\otimes\mathrm{id}_2,\ \mathrm{id}_1\otimes\varpi_2}"]\\
\&(E_1\otimes M_2)\times(M_1\otimes E_2)=\Pi\arrow[from=2-1,to=3-3,crossing over,"F_{12}" near end]\&
\end{tikzcd}}
\unsure{avec|car} 4 carrés commutatifs

Si on veut que $F$ induise \deleted{l'identité}\inserted{un hom d'extensions, donc} $L_1\otimes L_2\longrightarrow L_1\otimes L_2$, il f. et s. \deleted{\unsure{avoir}} que $\boxed{F_{22}=0}$, si on veut \inserted{de plus} que \unsure{ça|$\Sigma$} induise l'identité sur $L_1\otimes L_2$, il f. et s. que $\boxed{F_{21}=\alpha_1\otimes\alpha_2}$, \unsure{d\textsuperscript{o}|si|et} \unsure{si|{}} on veut que \unsure{ça|$F$} induise l'identité sur $M_1\otimes M_2$, il f. et s. que $\boxed{F_{12}=\psi_1\otimes\psi_2}$ : \deleted{\unsure{Si|D}} ces conditions sont vérifiées, alors \deleted{\unsure{d}}\unsure{deux} des 4 carrés voulus sont déjà commutatifs, savoir
\diagram{\begin{tikzcd}[ampersand replacement=\&,column sep=small,row sep=small]
\&\Sigma\arrow[dl,dash]\arrow[dr,dash]\&\\
E_1\otimes E_2\arrow[rrd,"0","{F_{12}=\psi_1\otimes\psi_2}"']\&\&L_1\otimes L_2\arrow[d,"{F_{22}=0}"]\\
\&\&M_1\otimes M_2
\end{tikzcd}}
et
\diagram{\begin{tikzcd}[ampersand replacement=\&,column sep=small,row sep=small]
\&\&L_1\otimes L_2\arrow[dll,"{\deleted{E_1\otimes E_2}F_{21}}"']\arrow[d,"{F_{22}=0}"]\\
E_1\otimes E_2\arrow[dr]\&\&M_1\otimes M_2\arrow[dl]\\
\&\deleted{\unsure{\varphi}}\Pi\&
\end{tikzcd}}
\margin{(les quatre composés sont nuls)}
% archive sheet number: 165
\scanpage{167}
\mspage{482}
re plus ou moins complet d'autres aspects encore de ma personne par lesquels se manifestait ce déséquilibre, dont j'ai essayé dans la note précédente de cerner aspect particulièrement apparent: celui du ``déployement de force''.

Il ne faudrait pas croire que ce déséquilibre, cultivé une vie durant, et la multitude des mécanismes psychiques par quoi il se manifestait, se soient évanouis du jour au lendemain comme par un coup de baguette magique. Je ne m'attendais à rien de tel d'ailleurs, ni en ce jour des retrouvailles, ni en les jours et semaines qui ont suivi.

(10 octobre) C'étaient des jours de fonte des glaces, portés par un afflux puissant d'énergie nouvelle - des jours de travail intérieur et d'émerveillement devant ces mondes nouveaux que jour après jour je voyais s'entrouvrir, prenant naissance dans l'humble trame des menus faits quotidiens et se déployant sous l'action intense d'yeux avides de voir. C'étaient les jours aussi où a commencé à poindre le premier pressentiment de la richesse de cet inconnu qui soudain m'interpellait, que j'avais ignoré la veille encore. Je l'appréhendais par ces ``bouts'' qui venaient de se faire connaître à moi, dans l'instant même des retrouvailles, et dans le voyage imprévisible et imprévu qui l'avait suivi. Je sentais bien que cette ``naissance'' par laquelle je venais de passer était tout juste le c\,o\,m\,m\,e\,n\,c\,e\,m\,e\,n\,t de quelque chose d'entièrement inconnu, ou plutôt le r\,e\,c\,o\,m\,m\,e\,n\,c\,e\,m\,e\,n\,t de quelque chose qui s'était interrompu, qui avait été coupé ou étouffé un jour, et qui était reparti mystérieusement. A vrai dire, ce ``d\handwritten{e}venir'' intense s'était remis en mouvement déjà dès les mois qui avaient précédé, mais à un niveau où la p\,e\,n\,s\,é\,e introspective n'avait guère eu encore de part...

Un des aspects profonds de ce devenir qui avait repris vie, de ce travail qui avait repris, a été la restauration progressive de l'équilibre originel de ``la femme'' et de ``l'homme'', d\handwritten{u} yin et du yang en moi, au fil des jours, des semaines et des années. D'une certaine façon, je peux dire que depuis le moment des retrouvailles, ``l'enfance'' ou l'état d'enfant est resté présent, ``en puissance'', par une connaissance profonde et indélébile en moi de ma propre nature, de mon u\,n\,i\,t\,é essentielle, indestructible, au delà des effets d'une certaine ``division'' qui souvent continue à agiter la surface de mon être. Le mot même ``enfant'' ou ``enfance'' pour désigner la c\,h\,o\,s\,e, cette unité de l'être, n'es\margin{\unsure{-}}
% archive sheet number: 166
\scanpage{168}
\diagram{\begin{tikzcd}[ampersand replacement=\&,row sep=small,column sep=tiny,cells={nodes={font=\scriptsize}}]
\&\&L_1\otimes L_2\arrow[dl]\arrow[dr]\&\&\&\&\&\\
\&E_1\otimes L_2\arrow[dr]\&\text{cocart}\&L_1\otimes E_2\arrow[dl]\&\&\&\&\\
0\arrow[rr]\&\&E_1\otimes_0 E_2\arrow[rr,"\mu"]\arrow[drr,phantom,"{\text{cocart}}" description]\arrow[d,"{\pi_1\otimes\mathrm{id},\mathrm{id}\otimes\pi_2}" description]\&\&E_1\otimes E_{\unsure{2|0}}\arrow[rr,dash]\arrow[d]\&\&M_1\otimes M_2\arrow[r]\arrow[d,"1"]\&0\\
0\arrow[rr]\&\&L_1\otimes L_2\arrow[rr]\arrow[dd,"1"]\&\&E_1\vee E_2\arrow[rr]\arrow[dd,"?","{\deleted{\unsure{F|E}}}"']\&\&M_1\otimes M_2\arrow[r]\arrow[dd,"1"]\&0\\
\rule{1.2cm}{.35pt}\&{}\&{}\&\rule{1.0cm}{.35pt}\&{}\&\rule{1cm}{.35pt}\&{}\&\leftarrow\mkern-10mu\rule{1.7cm}{.35pt}\\
0\arrow[rr]\&\&L_1\otimes L_2\arrow[rr]\arrow[d,"1"]\&\&E_1\wedge E_2\arrow[drr,phantom,"{\text{cart}}" description]\arrow[rr,dash]\arrow[d]\&\&M_1\otimes M_2\arrow[r]\arrow[d,"{\unsure{\varpi|\overline{\omega}}_1\otimes\mathrm{id},\ \mathrm{id}\otimes\varpi_2}"]\&0\\
0\arrow[rr]\&\&L_1\otimes L_2\arrow[rr]\&\&E_1\otimes E_2\arrow[rr]\&\&\deleted{M}E_1\otimes^{\unsure{0|\circ}}E_2\arrow[r]\arrow[dl]\arrow[dr]\&0\\
\&\&\&\&\&E_1\otimes M_2\arrow[dr,dash]\&\text{cart}\&\deleted{\unsure{L}}M_1\otimes E_2\arrow[dl,dash]\\
\&\&\&\&\&\&M_1\otimes M_2\&
\end{tikzcd}}
% archive sheet number: 167
\scanpage{169}
\mspage{483}
d'ailleurs apparu que des années plus tard, vers le moment où j'ai commencé à faire connaissance, au niveau de la pensée consciente, avec le double aspect yin-yang de toutes choses. C'était le moment aussi où apparaissait cette connaissance (ou du moins, ce pressentiment) que l'état d'enfance, l'état créateur, est celui de l'équilibre parfait des forces et énergies yin et yang, celui des ``épousailles'' du yin et du yang, se manifestant par un état d'harmonie créatrice.

Il me semble qu'à un certain niveau, cette connaissance de mon unité foncière est présente en tous les instants, et qu'elle a\,g\,i\,t en tout instant. Il est vrai aussi que cette action est plus ou moins sensible et efficace suivant les moments, et qu'elle n'est nullement dans la nature d'une élimination pl\unsure{us} ou moins permanente, voire d'une destruction en bloc des forces égotiques, du ``patron'' donc - ni même d'une élimination des forces de repression (qui forment une bonne part du ``moi'', sinon tout à fait sa totalité...). Ce sont là les force\unsure{s} de subreptice escamotage de la réalité qui m'entoure et de la réalité qui se dér\unsure{ou}le en moi - les forces silencieusement et obstinément à l'oeuvre pour maintenir contre vents et marées les tenaces illusions, qui sans elles s'effondreraient aussitôt sous leur propre poids... Certains de ces mécanismes de répression ont été repérés un à un et ont disparu. Je me suis débarassé de certaines i\,l\,l\,u\,s\,i\,o\,n\,s qui pesaient lourdement sur moi, et j'ai élucidé les quelques d\,o\,u\,t\,e\,s obstinés qui, pendant toute une vie, avaient été relégués (par les soins du ``patron'') croupir dans des souterrains-poubelle, jamais examinés. Leur message enfin entendu, ces doutes ont disparu, en laissant une connaissance paisible et joyeuse. J'ai repéré également des mécanismes de répression d'une grande puissance, profondément enracinés dans le moi, dont je réalise (depuis quelques années) que leur portée dans ma vie reste considérable aujourd'hui autant que jamais. Ils vont dans le sens du déséquilibre yang, dans le sens de l'occultation de certaines forces et facultés yin. J'ignore si ces mécanismes vont être désamorcés un jour - et je sais qu'il ne tient qu'à moi. Sans doute s'évanouiront-ils le jour, et le jour seulement, où je serai entré dans les origines du conflit \unsure{:|{}} dans ma vie bien plus profondément et plus totalement que je ne l'ai fait jusqu'à présent.

Pour le moment, avec l'orientation présente de ma vie vers un investissement mathématique important, je peux bien dire que ça n'en prend nullement le chemin!
% archive sheet number: 168
\scanpage{170}
\diagram{\begin{tikzcd}[ampersand replacement=\&,column sep=small]
0\arrow[r]\&L_1\arrow[r,"\alpha_1"]\&E_1\arrow[l,dashed,bend right=25,"\pi_1"']\arrow[r,"\psi_1"]\&M_1\arrow[l,dashed,bend right=25,"{\unsure{\varpi|\overline{\omega}}_1}"']\arrow[r,dash]\&0\\
0\arrow[r]\&L_2\arrow[r,dash,"\alpha_2"]\&E_2\arrow[l,dashed,bend left=25,"\pi_2"]\arrow[r,"\psi_2"]\&M_2\arrow[l,dashed,bend left=25,"\varpi_2"]\arrow[r,dash]\&0
\end{tikzcd}}
\margin{$\unsure{i|\iota|\alpha}_1,i_2$ \unsure{univ.} injectifs}

avec $\chi$-\unsure{trivialisation|trivialisations} $(\pi_1,\varpi_1)\;(\pi_2,\varpi_2)$
\[
\left\{\begin{aligned}
&\alpha_i\pi_i+\varpi_i\psi_i=\chi\,\mathrm{id}_{E_i}\\
&\deleted{\pi_1\varpi_2=\chi\,\mathrm{id}_{E_2}}\\
&\pi_i\,\deleted{\unsure{\vert L}}\alpha_i=\chi\,\mathrm{id}_{L_i}\quad
 \deleted{\unsure{\alpha_i\vert M_i/E_i}}=0\\
&\deleted{\unsure{\psi_i\varpi_i\vert L=0}}\ \psi_i\varpi_i
 \deleted{\unsure{+\chi\,\mathrm{id}_{M_i}}}\ \unsure{=}\chi\,\mathrm{id}_{M_{\unsure{i|2}}}
\end{aligned}\right.
\]
On définit de façon \unsure{duale|autre|évidente} $E_1\deleted{\unsure{\wedge}}\vee E_2$ et $E_1\wedge\deleted{L}E_2$, \unsure{extensions} $\unsure{\Pi|\pi}$-trivialisées de $M_1\otimes M_2$ par $L_1\otimes L_2$, et on a une application can. d'ext\deleted{\unsure{E}}ensions de $E_1\vee E_2$ \unsure{dans|sur|en} $E_1\wedge E_2$
\diagram{\begin{tikzcd}[ampersand replacement=\&,column sep=small,row sep=tiny]
\&E_1\otimes\deleted{\unsure{E}}L_2\arrow[dr,<->]\&\&\\
{[}L_1\otimes L_2\arrow[ur,dash]\arrow[dr]\&\&E_1\otimes_0 E_2\&\text{\unsure{cocartésien}}{]}\\
\&\deleted{\unsure{E}}L_{\unsure{1|2}}\otimes E_2\arrow[ur]\&\&
\end{tikzcd}}
\diagram{\begin{tikzcd}[ampersand replacement=\&,row sep=large,column sep=small]
\unsure{0}\arrow[r]\&E_1\otimes_0 E_2\arrow[r]\arrow[d,"{\pi_1\otimes\mathrm{id},\ \mathrm{id}\otimes\pi_2}"]\&E_1\otimes E_2\arrow[r]\arrow[d]\&M_1\otimes M_2\arrow[d,Rightarrow]\\
\unsure{0}\arrow[r,dash]\&L_1\otimes L_2\arrow[r]\&E_1\vee E_2\arrow[r]\&M_1\otimes M_2\\
0\arrow[r]\&L_1\otimes L_2\arrow[r]\arrow[d,Rightarrow]\&E_1\wedge E_2\arrow[r]\arrow[d]\&M_1\otimes M_2\arrow[d]\\
0\arrow[r]\&L_1\otimes L_2\arrow[r]\&E_1\otimes E_2\arrow[r]\&E_1\otimes^{\unsure{0|\circ}}E_2
\end{tikzcd}}
% archive sheet number: 169
\scanpage{171}
\mspage{484}
(\textsuperscript{110}) (11 octobre) J'ai eu envie depuis un jour ou deux de faire le point, en quelques mots, où en est (après huit ans) cette ``restauration progressive de l'équilibre yin-yang'' en moi.

Peut-être le changement le plus important de tous est dans une a\,c\,c\,e\,p\,t\,a\,t\,i\,o\,n beaucoup plus grande que par le passé de ma personne telle qu'elle est vraiment d'instant en instant. Une autre façon de l'exprimer, c'est que les mécanismes de répression en moi se sont considérablement assouplis. Comme je l'a dit hier, certains ont disparu après avoir été découverts et compris, et d'autres, que j'avais ignorés ma vie durant, me sont devenus familiers dans leurs manifestations de tous les jours. Je les vois en action, non comme des ennemis qu'il me faudrait essayer d'extirper coûte que coûte, mais comme faisant partie de la multiplicité des facettes de mon être conditionné, et par là, de la richesse du ``donné'' présent, lequel reflète fidèlement mon histoire passée; aussi bien l'histoire ``ancienne'' de mes conditionnements et des racines de la division dans mon être, que l'histoire plus récente de ma maturation, du travail donc par lequel je finis par déballer et par ``manger'' et assimiler le paquet initial légué par mes parents et par leurs successeurs. Cette ``acceptation'' en moi inclut donc: non seulement des pulsions et traits de ``l'enfant'' que j'avais pendant longtemps ignorés et réprimés (et notamment ceux qui reflètent les aspects féminins en moi), mais également les mécanismes de répression propres au ``patron'', c'est-à-dire justement des mécanismes invétérés de ``non-acceptation''! Accepter ces derniers n'a rien de commun avec ``les cultiver'', ou les fortifier. Au contraire, c'est un premier pas indispensable pour les dénouer ou les désamorcer tant soit peu, par l'effet d'une attention curieuse et aimante. L'expérience de ces huit années me donne la conviction que, pour peu que cette attention-là plonge assez profond et jusqu'à la racine même de la répression, celle-ci se résoud et disparaît en libérant une énergie considérable - celle qui jusque là était immobilisée pour maintenir contre vents et marées tel ensemble de mécanismes répressifs, et les habitudes de pensée et autres qui servent à les maintenir.

Mais ce n'est pas vis-à-vis des aspects par nature ``noués'' de ma personne, que cette acceptation nouvelle de moi-même a d'abord fait son apparition dans ma vie. Elle est venue sans tambour ni trompette, dès avant la découverte de la méditation, donc dès avant les ``retrouvailles'' la suivant de près. C'était au mois de juillet 1976, au cours d'une courte liaison amoureuse avec une jeune femme, G. , peut-être un brin plus ``homasse'' dans ses façons d'être que les femmes que j'avais aimées précédemment. Le hasard (?) a voulu que les circonstances matérielles qui ont entouré ces amours étaient telles, que je me voyais placé\margin{\unsure{-}} dans un rôle typiquement ``féminin''. Je faisais le ménage et préparais les
% archive sheet number: 170
\scanpage{172}
\[
u=u_1\overset{!}{\deleted{\unsure{\ast|\otimes|\star}}}u_2-b_1b_2
=2u_1u_2+b_2u_1+b_1u_2\ \in\ \unsure{A|\mathcal A}_1\otimes A_2
\]
solution de l'équation $\;x^2+bx+c=0$
\diagram{\begin{tikzcd}[ampersand replacement=\&,row sep=tiny]
A_1\otimes A_2\\ A\arrow[u,hook]
\end{tikzcd}}
(\unsure{injectif} sauf si car. 2 et $b_1=b_2=0$)
\[
\begin{aligned}
b&=b_1b_2\\
c&=b_1^2c_2+b_2^2c_1-4c_1c_2\\
\delta&=b^2-4c
=\underbrace{\unsure{\delta_1|\delta_1^2}}_{(b_1^2-4c_1)}
 \underbrace{\delta_2}_{(b_2^2-4c_2)}
\end{aligned}
\]
\diagram{\begin{tikzcd}[ampersand replacement=\&,column sep=large,row sep=large]
\&{\underline{0}\ \unsure{\mathcal O|O}}\&\\
A_1\otimes A_2\arrow[rr,bend left=25,"{\unsure{\leftrightarrows|\leftrightarrow|\sigma}}"]
\arrow[r,"{\substack{\text{can. par}\\\text{définition}\\\text{\unsure{directe|div.}}}}","{\text{de }A}"']\&
A_1\wedge \deleted{\unsure{A}}A_2=A\arrow[r]\&A_1\otimes A_2\\
E_1\otimes E_2\arrow[u,phantom,"{\unsure{\wedge|\uparrow}}" description]\arrow[d,dash]\&
\underset{\substack{\cup\\L_1\otimes L_2}}{E}=E_1\wedge E_2\arrow[l]\arrow[u,phantom,"{\unsure{\times|x}}" description]\arrow[d,dash]\&
E_1\otimes E_2\arrow[u,phantom,"{\unsure{\ast}}" description]\arrow[l,"{\text{can. par déf. \unsure{dér.|dir.|diff.}}}"',"{\text{de }E}"]\arrow[ll,bend left=25,"\leftrightarrows"]\arrow[d,dash]\\
0\&0\&0
\end{tikzcd}}
\[
\begin{aligned}
f\in A_1\otimes A_2\ \deleted{\unsure{\leftrightarrows(f)|\longmapsto}}\quad
f^{\leftrightarrows}\ \deleted{\unsure{=(\sigma_1\otimes\sigma_2+\mathrm{id})(f)}}&=
(\underbrace{\sigma_1\otimes\sigma_2}_{\sigma}+\mathrm{id}_{A_1\otimes A_2})(f)=f+\overline f\\
x\in E_1\otimes E_2\qquad x^{\leftrightarrows}&=x+\overline x
\end{aligned}
\]
\diagram{\begin{tikzcd}[ampersand replacement=\&,column sep=small,row sep=small]
\unsure{0}\arrow[d,dash]\&0\arrow[d,dash]\&0\arrow[d]\\
A_1\otimes_0 A_2\arrow[r,"\mathrm{Tr}"]\arrow[d]\&\mathcal O\arrow[r,Rightarrow,"\mathrm{id}"]\arrow[d,dash]\&A_1\otimes_0 A_2\arrow[d]\\
A_1\otimes A_2\arrow[r,dash]\arrow[d]\&A\arrow[r]\arrow[d,dash]\&A_1\otimes A_2\arrow[d]\\
L_1^\vee\otimes L_2^\vee\arrow[r,"\mathrm{id}"]\arrow[d]\&L_1^\vee\otimes L_2^\vee\arrow[r,"2"]\arrow[d,dash]\&L_1^\vee\otimes L_2^\vee\arrow[d]\\
0\&0\&0
\end{tikzcd}}
\margin{\deleted{$A\unsure{=|\longrightarrow}(A_1\otimes A_2)^{\leftrightarrows}$ est \unsure{dégénéré|définit} en car. 2}}
\diagram{\begin{tikzcd}[ampersand replacement=\&,column sep=small,row sep=small]
\unsure{0}\arrow[r,dash]\&L_1\arrow[r,dash]\&E_1\arrow[r]\&L'_1\arrow[r]\&0\\
0\arrow[r,dash]\&L_2\arrow[r,dash]\&E_2\arrow[r,dash]\&L'_2\arrow[r,dash]\&0
\end{tikzcd}}
\diagram{\begin{tikzcd}[ampersand replacement=\&,column sep=tiny,row sep=small]
\&E_1\otimes E_2\&\\
E_1\otimes\deleted{\unsure{L}}E_2\arrow[dr]\&\&L_1\otimes E_2\\
\&L_1\otimes L_2\&
\end{tikzcd}}
% archive sheet number: 171
\scanpage{173}
\mspage{491}
contact avec des choses décidément juteuses - chacun de ces groupes était riche de poids et de mystère, et chacun des couples yin-yang qui étaient censés le constituer (mais qui plutôt, tous ensemble, le d\,é\,s\,i\,g\,n\,e\,n\,t, sans aucunement l'épuiser) - chacun de ces couples a quelque chose de délicat et d'important à me dire sur la nature de ce monde dans lequel je vis, et souvent sur ma propre nature. J'ai retrouvé avec une force nouvelle ce sentiment qui était présent déjà il y a cinq ans: que le jeu délicat du yin et du yang, du ``féminin'' et du ``mâle'' en toutes choses, est un fil conducteur incomparable vers une compréhension du monde et de soi. Il nous conduit droit vers les questions essentielles. Souvent aussi, le ``yoga'' même du yin et du yang, le seul fait\handwritten{,} j'entends\handwritten{,} de faire attention à l'aspect des choses et évènements qui s'exprime en termes d'équilibre et de déséquilibre yin-yang, fournit une première clef pour une meilleure compréhension de ces questions, et vers une réponse.

Je m'excuse si pour certains lecteurs je dois donner l'impression, depuis une page ou deux, de parler du sexe des anges, alors qu'ils ne verraient pas trop même quels sont ces fameux ``couples'' yin- yang dont je parle, et encore moins ces ``groupes'' en lesquels certains se réunissent, lesquels groupes finalement seraient censés s'assembler en un ``diagramme'' (c'est quand même utile les maths!). Je devrais donner ici au moins un de ces groupes - et j'ai envie de prendre celui par lequel spontanément j'avais commencé hier, celui aussi qui a fini par apparaître au cours de la réflexion comme le groupe ``primitif'' (*), dont semblent sortir progressivement tous les autres, par des sortes de ``filiations'' successives (se poursuivant sur mon fameux diagramme sur huit ``générations''...). Voici donc la liste des ``couples'' que j'ai relevés, constituant ce groupe primitif (qu'on pourrait nommer par le premier de ces couples, savoir `` le groupe a\,c\,t\,i\,o\,n\handwritten{\unsure{$\checkmark$|,}}\,-\,i\,n\,a\,c\,t\,i\,o\,n'').

\begin{flushleft}
Action-inaction\\
activité-passivité\\
veille-sommeil\\
sujet-objet\\
engendrer-concevoir (*)\\
exécution-conception (*)
\end{flushleft}

\noindent\rule{3cm}{.35pt}

{\small (*) (6 novembre) En fait , il y a un groupe plus primitif encore, qu'on peut appeler le groupe ``p\,è\,r\,e\,-\,m\,è\,r\,e''. Voir, au sujet de cet ``oubli'', la note ``Notre Mère la Mort - ou l'Acte et le tabou'' (n° 113). Les couples ``engendrer-concevoir'' et\handwritten{``}exécution-conception'', que j'ai fait figurer ci-dessous dans le groupe (soi-disant ``primitif'') \handwritten{ac}tion-inaction, s'insèrent visiblement de façon plus naturelle dans le ``groupe mère'' formé autour du couple ``père-mère''.}
% archive sheet number: 172
\scanpage{174}
\diagram{\begin{tikzcd}[ampersand replacement=\&,column sep=small]
\unsure{0}\arrow[r,dash]\&L\arrow[r,"i"]\&E\arrow[r,"{\unsure{\psi|\varphi}}"]\&L'\arrow[r,dash]\&0
\end{tikzcd}}
\[
\begin{aligned}
E\xrightarrow{\pi}L&\qquad\text{tel que}\qquad \pi(x)=\chi x\quad\text{si }x\in L\\
L'\xrightarrow{\unsure{\varpi|\overline{\omega}}}E&\qquad\text{tel que}\qquad
\psi(\deleted{\unsure{\pi}}\varpi(\lambda))=\chi\lambda\quad\text{si }\lambda\in L\\
E\xrightarrow{\deleted{\unsure{\phi}}f}E&\qquad\text{tel que}\qquad
1^\circ)\ f\ \text{nul sur } L\\
&\hspace{7em}2^\circ)\ \deleted{f(x)-\unsure{fx|\chi x}\in L}\\
&\hspace{8em}\underbrace{\chi x-f(x)}_{\substack{\deleted{f-\chi\,\mathrm{id}}\\\pi(x)}}\in L
\end{aligned}
\]
\margin{donc \unsure{se factorise} en $E\xrightarrow{\psi}L'\xrightarrow{\varpi}E$}

\[
\boxed{\pi=\chi\,\mathrm{id}-f}
\]
\[
\begin{aligned}
&\deleted{\pi:\text{\unsure{valeurs|identité} dans }L\ \longleftrightarrow}\\
&\unsure{\varpi|\pi}\ \text{stabilise }L\ \longleftrightarrow\ f\ \text{stabilise }L\\
&\pi:\text{ val. dans }L\ \inserted{\text{i.e. nul sur }L'=E/L}
 \longleftrightarrow f\ \text{induit }\deleted{\text{l'identité}}\chi\,\mathrm{id}_{L'}\ \text{sur }E/L\\
&\pi\ \text{induit }\chi\,\mathrm{id}_L\ \text{sur }L\ \longleftarrow f\ \text{nul sur }L'
\end{aligned}
\]
\diagram{\begin{tikzcd}[ampersand replacement=\&,row sep=small,column sep=small]
0\arrow[r]\&{L'}^\vee\arrow[r,<->]\&\unsure{A|\mathcal A}\arrow[r,Rightarrow]\&L^\vee\arrow[r,dash]\&0\\
0\arrow[r]\&L\arrow[r]\&E\arrow[r,"{\unsure{\mapsto|{}|n}}"']\&L'\arrow[r,dash]\&0
\end{tikzcd}}
\diagram{\begin{tikzcd}[ampersand replacement=\&,column sep=small]
\deleted{0}\arrow[rrrr,dash,start anchor=west,end anchor=east]\arrow[r,dash]\&\deleted{L\otimes M}\arrow[r]\&\deleted{E\otimes M}\arrow[r,dash]\&\deleted{L'\otimes M}\arrow[r]\&\deleted{0}
\end{tikzcd}}
\[
\deleted{E_1\unsure{\wedge|\otimes}E_2}
\]
\diagram{\begin{tikzcd}[ampersand replacement=\&,row sep=small,column sep=small]
0\arrow[r]\&L_1\arrow[r]\&E_1\arrow[r]\&L'_1\arrow[r]\&0\\
0\arrow[r,dash]\&L_2\arrow[r,dash]\&E_2\arrow[r,dash]\&L'_2\arrow[r,dash]\&0
\end{tikzcd}}
\diagram{\begin{tikzcd}[ampersand replacement=\&,row sep=small,column sep=small]
\unsure{0}\arrow[r]\&L_{\unsure{1|2}}\arrow[r,dash]\&E_{\unsure{1|11}}\arrow[r,dash]\&L'_1\arrow[r,dash]\&0\\
0\arrow[r]\&L_2^{\unsure{\prime|{}}}\arrow[r]\&E_2\arrow[r,dash]\&L'_2\arrow[r]\&0
\end{tikzcd}}
\[
E_1\simeq A_1\otimes\Delta_1
\qquad E_1\wedge E_2
\]
\[
\begin{aligned}
\overset{\Delta_1^{-1}}{E_1^\vee}\wedge\overset{\Delta_2^{-1}}{E_2^\vee}
&\simeq\underbrace{(E_1\wedge E_2)}_{\Delta_1\ \Delta_2}
 \otimes\Delta_1^{\unsure{\prime|-1}}\otimes\Delta_2^{-1}\\
&\simeq(E_1\wedge E_2)^\vee
\end{aligned}
\]
\[
\begin{aligned}
\pi(x)&=2x-\deleted{f}F(x)\\
\pi(x)-x&=x-F(x)
\end{aligned}
\]
% archive sheet number: 173
\scanpage{175}
\mspage{490}
jours. Puis, entraîné par le petit diable (ou ange, je ne sais...) de la systématique en moi, j'ai fini par ressortir mes \handwritten{vieilles réflexions d'il y a} cinq ans sur ce thème. Pendant une semaine ou deux je m'étais amusé alors à ``ramasser'' une centaine ou deux de ces couples bien suggestifs, lesquels s'étaient alors assemblés par affinités en une vingtaine de groupes. Comme cette réflexion s'est faite en marge du fameux ``ouvrage poétique'' que j'étais en train d'écrire, je n'ai pu m'empêcher de ranger ces groupes tant bien que mal à la queue-leu-leu, par affinités et filiations de sens d'un groupe \handwritten{au sui}vant. Hier soir, reprenant la réflexion avec du recul, et sans carcan poétique autour du cou, c'est dix-huit groupes que j'ai trouvé\handwritten{s} (au lieu de vingt), par un groupage peut-être un peu plus rigoureux. Je soupçonne d'ailleurs qu'il doit y avoir bien d'autres groupes encore, peut-être même un nombre illimité, correspondants à des modes d'appréhension de la réalité auxquels je n'ai pas songé au cours du travail (ni, peut-être, jamais encore).

Quant aux dix-huit groupes que j'ai bel et bien relevés, je me suis efforcé de les assembler en un diagramme (ou ``graphe'') suivant les principaux liens d'affinités qui les relient les uns aux autres. Certains de ces liens d'ailleurs ne se sont imposés à mon attention qu'au cours du tracé d'ébauches successives du diagramme. Le travail ici était vraiment très proche du travail mathématique bien familier, quand on s'efforce de saisir graphiquement, de façon aussi frappante que possible, un ensemble plus ou moins complexe de relations (données par exemple par des ``aplications'', figurées par des flèches) entre un certain nombre d'``ensembles'' ou de ``catégories'', figurant comme ``sommets'' du ``diagramme'' qu'on s'efforce de construire. Là aussi, des exigences de nature essentiellement esthétiques, de symétrie et de transparence structurale notamment, conduisent fréquemment à introduire (et au besoin donc, à découvrir voire même à inventer) des ``flèches'' ou liens auxquels on n'avait pas songé au départ, et parfois même des nouveaux ``sommets''. Toujours est-il qu'après cinq ou six ébauches successives, j'ai fini par aboutir à un diagramme, vaguement en forme d'arbre de Noël, qui m'a satisfait provisoirement - d'autant plus qu'il commençait vraiment à se faire prohibitivement tard!

Je me suis couché content, je sentais que je n'avais pas perdu mon temps, même si nes notes n'avaient pas avancé d'un poil (*). Mais je m'étais remis en

\noindent\rule{3cm}{.35pt}

{\small (*) En compensation, je pourrais déposer un brevet sur l'invention d'une nouvelle forme \handwritten{poétique, savoir le poème dit ``non linéaire''; ou ``diagrammatique''.}}
% archive sheet number: 174
\scanpage{176}
\diagram{\begin{tikzcd}[ampersand replacement=\&,column sep=tiny,row sep=small]
\unsure{A|\mathcal A}_{i}\&\&\&A\&u\&1\\
E_i\&x_i,\& e_i\&\deleted{\unsure{e}}E\&x_1\unsure{\ast|\otimes|\star}x_2=x,\&e=e_1\otimes e_2\\
\&P_{\unsure{1|i}}\arrow[u]\&L_1\arrow[u]\&\&P_{\unsure{3|2}}\arrow[u]\&
\end{tikzcd}}
\[
\begin{aligned}
\langle u_i,\unsure{x_i|u_i}\rangle&=\langle1_{\unsure{i|j}},\unsure{e_i|e'}\rangle=0&
 \langle u,e\rangle&=\langle1,x\rangle=1\\
\langle u_i,e_i\rangle&=\langle1_i,x_i\rangle=1&
 \langle u,x\rangle&=\langle1,e\rangle=0
\end{aligned}
\]
\[
\begin{aligned}
u_1&\in A_1&\deleted{\unsure{A\longrightarrow A_1\otimes A_2|A_1\otimes A_2}}\\
u_2&\in A_2&
\end{aligned}
\]
On veut exprimer
\[
\deleted{\alpha(u_1\ast u_2)=2u_1\unsure{u_2}+b_2u_1+}
\]
$u_1\ast u_2$ comme $\lambda\cdot1_A+\mu u$, \unsure{donc|en|dans} \unsure{calculer|calculs}
\[
\begin{aligned}
\lambda&=\underbrace{\langle u_1\ast u_2,x\rangle}_{\unsure{\Vert|u}} ,
 &\mu&=\langle u_1\ast u_2,e\rangle
\end{aligned}
\]
\[
\begin{aligned}
\deleted{A_1\otimes A_2\longrightarrow A}\quad
\langle u_1\ast u_2,x_1\ast x_2\rangle
 &=\langle\underbrace{u_1\otimes u_2+\overline u_1\otimes\overline u_2}_{(u_1\otimes u_2)^{\unsure{\leftrightarrows|\leftrightarrow|\sigma}}},x_1\otimes x_2\rangle\\
 &\deleted{\unsure{=u_1}}\ \deleted{\unsure{\langle x}}\langle 2u_1\deleted{\unsure{\otimes}}u_2+
 \underset{!}{b_1}u_2+u_1b_2+b_1b_2,\ x_1\otimes x_2\unsure{\rangle}\\
 &\deleted{=2\unsure{+\beta_1+\beta|\beta_1+\beta}\quad\unsure{=\beta_1}}\quad b_1b_2
\end{aligned}
\]
\[
\begin{aligned}
&\deleted{u_1\ast u_2,}\\
\mu&=\langle u_1\ast u_2,e\rangle
 =\deleted{\langle 2u_1\otimes u_2+b_1u_2+u_1b_2+b_1b_2,e_1\otimes e_2\rangle}\\
&\quad \deleted{\unsure{\varphi}}\inserted{\langle\varphi_1(u_1),e_1\rangle
 \langle\varphi_2(u_2),e_2\rangle}\\
&\deleted{=2\unsure{,}}\qquad=1\qquad e\in\Gamma L\subset\Gamma\deleted{\unsure{A}}E
\end{aligned}
\]
donc
\[
u_1\ast u_2=b_1b_2+\deleted{\unsure{2}}u
\qquad\deleted{\unsure{u=}}
\]
i.e.
\[
\boxed{u_1\ast u_2-b_1b_2=\deleted{\unsure{2}}u}
\]
\[
u_1\ast u_2-\mathrm{Tr}_1(u_1)\mathrm{Tr}_2(u_2)=2\cdot
\]
l'image de $u$ dans $A_1\otimes A_2$ est donc
\[
\boxed{2u_1\otimes u_2+b_2u_1+b_1u_2\unsure{+|{}}}
\]
% archive sheet number: 175
\scanpage{177}
\mspage{492}

\begin{flushleft}
dynamisme-équilibre\\
élan-assise\\
ardeur-persévérance\\
fougue- patience\\
passion-sérénité\\
tenacité-détachement.
\end{flushleft}

J'y joindrais bien encore les deux couples suivants, parmi une dizaine de ``retardataires'' qui me sont venus encore ce matin, sur la lancée de ma réfléxion de hier:

\begin{flushleft}
savoir-connaître\\
expliquer-comprendre.
\end{flushleft}

Est-il besoin de préciser que dans ces couples, c'est le terme ``yang'' ou ``masculin'' qui est mis le premier, suivant l'usage de notre société patriarcale, où l'homme donne le nom au couple? Par contre, alors que la société chinoise traditionelle est considérablement plus patriarcale encore que la nôtre, quand on suit l'usage chinois pour parler de la relation du yin et du yang, on met toujours en premier le yin (``féminin''), en parlant p. ex. d'``équilibre yin-yang'' (au lieu de yang-yin). Le sens de cet usage est sûrement dans l'intuition-archétype que c'est le yang qui naît du yin, lequel est le principe ``le plus primitif'' des deux, et non l'inverse...

Ce n'est pas le lieu ici de me lancer dans des comentaires sur l'un ou l'autre de ces couples. Pour le lecteur qui ne ``sentirait rien'' en les voyant, ce serait de toutes façons peine perdue; et celui qui se sent interpellé par eux, qui sent (fut-ce obscurément) que chacun d'eux a quelque chose à lui dire sur le monde et sur lui-même - sur l'équilibre et le déséquilibre, sur la dynamique interne des êtres et des choses..., celui-là peut se passer de commentaires circonstanciés, et prendre cette interpellation comme un point de départ pour sa propre réflexion.

\handwritten{(\unsure{111'|111$^{1}$})} Il y a un point seulement sur lequel je voudrais insister ici, commun à tous les ``couples'' yin-yang sans exception. C'est la chose aussi la plus cruciale de toutes, il me semble, pour une compréhension de la nature de la relation entre le yin et le yang, et par là,\handwritten{\deleted{\illegible{1}}} de la nature de chacun de ces deux principes (ou énergies, ou aspects, ou forces...) dans l'Univers. C'est ceci: chacun des deux termes d'un de ces couples, tel action-inaction, e\,n\quad l\,'\,a\,b-
\margin{\unsure{-}}
% archive sheet number: 176
